\documentclass[11pt,a4paper,DIV=12]{scrartcl}
\usepackage{multicol}
\makeatletter
\def\input@path{{../../../}}
\makeatother

\usepackage[sexy]{evan}
% \usepackage[sexy]{../../../evan}

% Override environment "problem" agar hanya menampilkan angka
\makeatletter
\newtheoremstyle{justnumber}
  {8pt} % space above
  {8pt} % space below
  {\normalfont} % body font
  {} % indent amount
  {\color{blue!40!black}\bfseries} % head font
  {} % punctuation after head
  {0.5em} % space after head
  {\thmnumber{#2}.\thmnote{ (\small #3)}} % head spec (hanya angka dan note)

% Hapus definisi problem lama di dalam makeatletter agar @ terbaca
\let\problem\relax
\let\endproblem\relax
\let\c@problem\relax
\let\theproblem\relax
\makeatother

\newcommand\answerbox{%%
\vspace*{0.5cm}%
\noindent\textbf{\textsf{Jawab}}: \fbox{\rule{14cm}{0pt}\rule[-0.2ex]{0pt}{1.9ex}}}

\theoremstyle{justnumber}
\newtheorem{problem}{}

\begin{document}
\title{Tryout 1 Pembinaan OSK Matematika}
\author{MAN 1 Pasuruan}
\date{23 April 2025}
\maketitle
\noindent\textbf{Petunjuk pengerjaan:}
\begin{itemize}
  \item Waktu pengerjaan adalah 120 menit.
  \item Terdapat 20 soal yang terdiri dari 10 soal pilihan ganda dan 10 soal isian singkat.
  \item Nilai soal pilihan ganda adalah {\color{green!70!black}$B=+1$}, {\color{red}$S=0$}, $K=0$.
  \item Nilai soal isian singkat adalah {\color{green!70!black}$B=+4$}, {\color{red}$S=0$}, $K=0$.
  \item Gunakan bolpoin untuk menulis jawaban, dan coretlah jawaban yang salah.
\end{itemize}

\section*{Pilihan Ganda}
\begin{problem}
  Tentukan banyaknya bilangan ganjil dari barisan
  $$0,1,2,3,4,5,6,7,8,9,1,0,1,1,\ldots,2,0,2,5,2,0,2,6$$
  \begin{multicols}{5}
    \begin{enumerate}
      \item[(a)] 4020
      \item[(b)] 4021
      \item[(c)] 4022
      \item[(d)] 4023
      \item[(e)] 4024
    \end{enumerate}
  \end{multicols}
  %KJ: D
\end{problem}

\begin{problem}
  Suatu bilangan asli tiga digit $\overline{abc}$ disebut sebagai bilangan \textit{imoetz} jika digit tengah $b$ merupakan kelipatan dari selisih $a-c$. Tentukan banyaknya bilangan \textit{imoetz} yang memenuhi kondisi tersebut.
  \begin{multicols}{5}
    \begin{enumerate}
      \item[(a)] 386
      \item[(b)] 387
      \item[(c)] 388
      \item[(d)] 389
      \item[(e)] 390
    \end{enumerate}
  \end{multicols}
  %KJ: D
\end{problem}

\begin{problem}
  Ucup sedang menyusun kata-kata yang terdiri dari huruf-huruf
  $$M,E,J,I,K,U$$
  Tentukan banyaknya susunan kata yang terdiri dari 5 huruf, dengan syarat bahwa jika huruf $J$ dan $K$ digunakan dalam susunan tersebut, maka tidak boleh bersebelahan.

  \begin{multicols}{5}
    \begin{enumerate}
      \item[(A)] 720
      \item[(B)] 624
      \item[(C)] 528
      \item[(D)] 288
      \item[(E)] 144
    \end{enumerate}
  \end{multicols}
  %KJ: C
\end{problem}

\begin{problem}
  Enam kartu yang diberi nomor 1 sampai 6 akan disusun berjajar dalam satu baris. Tentukan banyaknya susunan kartu tersebut sehingga terdapat tepat satu kartu yang dapat dihapus, sehingga lima kartu yang tersisa membentuk urutan menaik (ascending) atau menurun (descending).
  \begin{multicols}{5}
    \begin{enumerate}
      \item[(A)] $60$
      \item[(B)] $58$
      \item[(C)] $56$
      \item[(D)] $54$
      \item[(E)] $52$
    \end{enumerate}
  \end{multicols}
  %KJ: E
\end{problem}

\begin{problem}
  Suatu pesta dihadiri oleh remaja (laki-laki dan perempuan), laki-laki dewasa, dan perempuan dewasa. Dalam pesta tersebut, remaja hanya berjabat tangan sekali dengan sesama remaja, laki-laki dewasa yang hanya berjabat tangan sekali dengan sesama laki-laki dewasa, begitu pula dengan perempuan dewasa. Jika banyaknya laki-laki dewasa yang hadir dalam pesta lebih banyak dari perempuan dewasa dan remaja serta jumlah seluruh jabat tangan adalah $21$ jabat tangan. Banyaknya laki-laki yang hadir dalam pesta tersebut adalah \ldots
  \begin{multicols}{5}
    \begin{enumerate}
      \item[(A)] $1$
      \item[(B)] $3$
      \item[(C)] $5$
      \item[(D)] $6$
      \item[(E)] $7$
    \end{enumerate}
  \end{multicols}
  %KJ: D
\end{problem}

% \begin{problem}

%   \begin{multicols}{5}
%     \begin{enumerate}
%       \item[(A)] $-2$
%       \item[(B)] $\dfrac{1}{2}$
%       \item[(C)] $2$
%       \item[(D)] $\dfrac{1}{4}$
%       \item[(E)] $4$
%     \end{enumerate}
%   \end{multicols}
%   %KJ: C
% \end{problem}



\section*{Isian Singkat}
\begin{problem}
  Ada 10 orang, lima laki-laki dan lima perempuan, termasuk sepasang
  pengantin. Seorang tukang foto yang bukan salah satu di antara 10 orang
  tersebut akan megambil gambar enam orang di antara mereka, termasuk kedua pengantin, dengan tidak ada laki-laki dan perempuann yang
  berdekatan kecuali pengantin. Banyaknya cara adalah \ldots

  %KJ: 768
  \answerbox
\end{problem}

\begin{problem}
  Bola bernomor $1, 2, 3, \ldots$ dimasukkan ke dalam 5 kotak yang diberi label $A, B, C, D,$ dan $E$ dengan prosedur berikut. Bola 1 dimasukkan ke kotak $A$, dan bola 2 serta 3 dimasukkan ke kotak $B$. Tiga bola berikutnya dimasukkan ke kotak $C$, empat bola berikutnya ke kotak $D$, dan seterusnya, kembali ke kotak $A$ setelah bola dimasukkan ke kotak $E$. (Sebagai contoh, bola 22, 23, \ldots, 28 dimasukkan ke kotak $B$ pada langkah ke-7 dari proses ini.) Pada kotak manakah bola 2026 dimasukkan?

  %KJ: D
  \answerbox
\end{problem}

\begin{problem}
  Diberikan suatu dadu tidak standar dengan bilangan pada sisi-sisinya $3, 5, 8, 13, 21,$ dan $34$. Dadu tersebut dilemparkan dua kali. Banyaknya kemungkinan jumlah bilangan yang muncul merupakan suatu bilangan pada sisi dadu tersebut adalah \ldots

  %KJ: 8
  \answerbox
\end{problem}

\begin{problem}
  Tentukan banyaknya garis berbeda pada bidang Kartesius yang dapat dibentuk dari pasangan $(a,b) \in \{0,1,2,3,6,7\}$ sehingga persamaan garis tersebut adalah
  $ax + by = 0$.\\
  \textbf{Hint:} Gunakan gradien garis ketika $a,b \neq 0$.

  %KJ: 19
  \answerbox
\end{problem}

\begin{problem}
  Ucup menulis bilangan $26$ pada papan tulis. Kemudian ia menyisipkan $3$ angka berbeda di depan, di tengah, atau di belakang sehingga terbentuk bilangan dengan $5$ digit berbeda.
  Banyak bilangan yang dapat ia peroleh adalah \ldots

  %KJ: 3108
  \answerbox
\end{problem}

\begin{problem}
  Carilah nilai dari $A/1012$, dimana
  $$A = \frac{1^3 + 1}{1 + 1} + \frac{2^3 + 1}{2 + 1} + \frac{3^3 + 1}{3 + 1} + \cdots + \frac{2024^3 + 1}{2024 + 1}$$

  %KJ: 2731052 
  \answerbox
\end{problem}

\begin{problem}
  Diketahui $f(x)$ adalah fungsi yang memenuhi
  $$f(x-y) = f(x^2-y^2)$$
  untuk semua bilangan bulat $x, y$. Jika $f(2024) = 2025$ dan $f(2025) = -2024$. Maka carilah nilai dari $f(1) + f(2) + \dots + f(2024)$

  %KJ: 1012
  \answerbox
\end{problem}

\end{document}